\section*{Hoofdstuk \ref{ch:b2:wilson-cycle} --- Oude massieven en de Wilsoncyclus}

\begin{solution}{exo:b2:wilson-cycle:1}
(a) $600 - 400 = \qty{200}{Myr}$, van het uiteenscheuren van de randen tot de laatste botsing.
(b) $480 - 400 = \qty{80}{Myr}$. (c) De Caledonische orogenese is een reeks fasen, verspreid over
zo'n 200 miljoen jaar, en elke fase begon op verschillende plaatsen op verschillende tijdstippen: het woord
situeert een oceaan en zijn gordel, geen moment. Een ouderdom wordt gegeven door de stratigrafische schaal of in
miljoenen jaren.
\end{solution}

\begin{solution}{exo:b2:wilson-cycle:2}
(c) de Oost-Afrikaanse rift, riftvorming; (a) de Rode Zee, een jonge oceaan; (d) de Atlantische Oceaan, een wijde
oceaan met passieve randen; (e) de Pacifische rand van Zuid-Amerika, subductie en het
sluiten van een oceaan; (b) de Himalaya, botsing; (f) Bretagne, een afgevlakte gordel, waarna
de riftvorming opnieuw kan beginnen.
\end{solution}

\begin{solution}{exo:b2:wilson-cycle:3}
(a) Na de afzetting van de vroeg-Silurische lagen, dus na \qty{443.1}{Ma}, en vóór de
afzetting van de lagen uit het Laat-Devoon, dus vóór \qty{358.9}{Ma}, het einde van het Devoon.
(b) Minder dan $443.1 - 358.9 = \qty{84.2}{Myr}$; in werkelijkheid minder, omdat de onderste lagen jonger zijn
dan de basis van het Siluur en de bovenste lagen ouder dan de top van het Devoon. (c) Afzetting van de
grijze turbidieten in de oceaan; plooiing en splijting in de Caledonische botsingen; opheffing en
erosie tot een oneffen landoppervlak; afzetting van de rode breccies en conglomeraten; zachte
kanteling in de Varistische orogenese; de huidige erosie, die de klif uitsnijdt.
\end{solution}

\begin{solution}{exo:b2:wilson-cycle:4}
(a) Beide: een blok continentale korst dat door oceanen van de continenten gescheiden was terwijl het
de Iapetus overstak, en, eenmaal aangemeerd, een door breuken begrensd blok waarvan de geschiedenis verschilt van die van zijn
buren. (b) Geen van beide: een continent, geen blok dat aan een continent is toegevoegd. (c) Een \omterm{def:b2:wilson-cycle:terrane}{terrane}, geen
\omterm{def:b2:wilson-cycle:terrane}{microcontinent}: zijn korst is die van een boog, geen fragment van een continent. (d) Geen van beide: het is
de grens tussen \omterm{def:b2:wilson-cycle:terrane}{terranes}, het spoor van de gesloten oceaan.
\end{solution}

\begin{solution}{exo:b2:wilson-cycle:5}
(a) $3000 - 2600 = 400$, $2600 - 1500 = 1100$, $1500 - 875 = 625$, $875 - 600 = 275$, $600 -
262.5 = 337.5$; gemiddelde $2737.5/5 = \qty{547.5}{Myr}$. (b) $1100/275 = 4.0$. (c) $(1500 -
262.5)/3 = \qty{412.5}{Myr}$. Het gemiddelde verschuift met een kwart wanneer de twijfelachtige punten worden weggelaten,
en de tussenafstanden variëren met een factor vier: een ruwe herhaling, geen periode.
\end{solution}

\begin{solution}{exo:b2:wilson-cycle:6}
(a) $3000/(70 - 25) = \qty{66.7}{Myr}$. (b) $3000/70 = \qty{42.9}{Myr}$. (c) Zodra de rug is
gesubduceerd, maakt de overblijvende oceaanbodem deel uit van de verre plaat, die het verre
continent draagt: de sluitingssnelheid is de snelheid van die plaat ten opzichte van de trog.
\end{solution}

\begin{solution}{exo:b2:wilson-cycle:7}
(a) $700/8 = \qty{87.5}{km}$; $700/30 = \qty{23.3}{km}$. (b) De Varistische gordel: bij
\qty{30}{\celsius} per kilometer smelt de middenkorst; bij \qty{8}{\celsius} per kilometer
ligt de smelttemperatuur onder de basis van zelfs een verdikte korst. (c) De lage gradiënt
beschrijft gesteenten die snel in een subductiezone omlaag werden gevoerd, voordat ze konden opwarmen. Gesteenten die
diep in een verdikte korst blijven, warmen over miljoenen jaren op naar een normale geotherm, en kunnen
dan op kleinere diepte het smelten bereiken.
\end{solution}

\begin{solution}{exo:b2:wilson-cycle:8}
(a) $\tau\ln(3/0.3) = 26.3\times2.303 = \qty{60.6}{Myr}$. (b) $290 - 60.6 = \qty{229}{Ma}$, in
het Laat-Trias. (c) $\tau = (290 - 251.9)/\ln10 = \qty{16.5}{Myr}$. (d) De extensionele
ineenstorting van de verdikte korst, die haar dunner maakt en haar oppervlak rechtstreeks verlaagt; ook een
hogere erosiesnelheid, een natter klimaat of zachtere gesteenten, of een lagere beginhoogte.
\end{solution}

\begin{solution}{exo:b2:wilson-cycle:9}
(a) De sutuur, het spoor van de Iapetus, ligt op het land in Groot-Brittannië en Ierland, ten oosten van de
Atlantische Oceaan; korst van Laurentia, van de overkant van de Iapetus, ligt bij Rockall aan de
Europese kant. De nieuwe rift sneed ten westen van de oude sutuur: hij ging ernaast open, niet erlangs.
(b) Ten noorden van de sutuur sokkel en dek van de Laurentische rand; ten zuiden ervan die van
Avalonia, met verschillende geschiedenissen en, in het vroege Paleozoïcum, verschillende fauna's. (c) De riftvorming
kwam na 80 miljoen jaar, meer dan de 30 tot 50 waarna een afgekoelde zone sterk kan zijn
geworden; toch volgde de Centrale Atlantische Oceaan wel de Appalachen. De warmte van een gordel vervaagt; zijn
breuken, maaksels en samenstellingen blijven en kunnen nog steeds een rift leiden.
\end{solution}

\begin{solution}{exo:b2:wilson-cycle:10}
(a) $70 - 40/2 = \qty{50}{mm}$ per jaar. (b) $\qty{1500}{km}/\qty{50}{km.Myr^{-1}} =
\qty{30}{Myr}$. (c) De oceaan versmalt met $70 - 40 = \qty{30}{mm}$ per jaar: $3000 - 30\times30 =
\qty{2100}{km}$. (d) Jonge, hete lithosfeer komt dan bij de trog aan, en er dehydrateert geen oude, koude plaat
meer onder de boog: het boogmagmatisme verandert of stopt.
\end{solution}

\begin{solution}{exo:b2:wilson-cycle:11}
(a) $h = \Delta T(\rho_m - \rho_c)/\rho_m = 25\times600/3300 = \qty{4.5}{km}$. (b) De hele extra
\qty{25}{km}, 5.5 keer de hoogte van de keten. (c) $2700\times9.8\times25\,000 =
\qty{6.6e8}{Pa}$, ongeveer \qty{0.66}{GPa}. (d) Bij ongeveer \qty{0.7}{GPa} en de temperaturen van
de wortel van een verdikte korst, enkele honderden graden, de amfiboliet- of de granulietfacies:
de gneisen en \omterm{def:b2:crustal-melting:migmatite}{migmatieten} van de oude massieven.
\end{solution}

\begin{solution}{exo:b2:wilson-cycle:12}
(a) $S = 30\times10 = \qty{300}{km}$; $W_0 = 300 + 300 = \qty{600}{km}$; $T = 35\times600/300 =
\qty{70}{km}$. (b) $h = 35\times600/3300 = \qty{6.36}{km}$; $r = 4.5h = \qty{28.6}{km}$.
(c) $\Delta F = 2700\times9.8\times6364\times(35\,000 + 3182 + 14\,318) =
\qty{8.8e12}{N.m^{-1}}$, 1.4 keer die van Tibet. (d) Zo'n gordel is zwaarder dan Tibet, dat al
uitvloeit: \omterm{def:b2:himalaya-andes:collapse}{gravitationele ineenstorting}, door rek van zijn hoge gebieden en overschuiving aan zijn randen.
\end{solution}

\begin{solution}{pb:b2:wilson-cycle:1}
\textbf{1.} $720/(540 - 500) = \qty{18}{km}$ per miljoen jaar, \qty{18}{mm} per jaar.
\textbf{2.} \qty{36}{mm} per jaar: een traag spreidende rug (20 tot \qty{55}{mm} per jaar).
\textbf{3.} $36\times(580 - 540) = \qty{1440}{km}$.
\textbf{4.} $36\times(580 - 500) = \qty{2880}{km}$.
\textbf{5.} Symmetrische spreiding met constante snelheid vanaf breedte nul op \qty{580}{Ma}, en geen
subductie vóór \qty{500}{Ma}.
\textbf{6.} \qty{80}{Myr}, aan beide randen van de oceaan, tegen de continentale randen. Het is de
koudste, dikste en dichtste oceanische lithosfeer van de oceaan, die het meest geneigd is te zinken.
\textbf{7.} $64 - 36 = \qty{28}{mm}$ per jaar.
\textbf{8.} $2880/28 = \qty{102.9}{Myr}$ na \qty{500}{Ma}: op \qty{397.1}{Ma}.
\textbf{9.} De rug, \qty{1440}{km} van de trog, nadert die met $64 - 18 =
\qty{46}{mm}$ per jaar: na $1440/46 = \qty{31.3}{Myr}$, op \qty{468.7}{Ma}.
\textbf{10.} $2880 - 28\times31.3 = \qty{2004}{km}$.
\textbf{11.} $2004/40 = \qty{50.1}{Myr}$ later: op $468.7 - 50.1 = \qty{418.6}{Ma}$.
\textbf{12.} Binnen de 425 tot \qty{400}{Ma} van de Scandische botsing.
\textbf{13.} $40\times8 = \qty{320}{km}$.
\textbf{14.} $W_0 = 400 + 320 = \qty{720}{km}$; $T = 35\times720/400 = \qty{63.0}{km}$.
\textbf{15.} $h = 28\times600/3300 = \qty{5.09}{km}$; $r = 4.5\times5.09 = \qty{22.9}{km}$.
\textbf{16.} $\Delta F = 2700\times9.8\times5091\times(35\,000 + 2545 + 11\,455) =
\qty{6.60e12}{N.m^{-1}}$, ongeveer \qty{8}{\percent} meer dan die van Tibet.
\textbf{17.} Stoppen met rijzen en uitvloeien: rek en afschuivingen in zijn hoge binnenste,
overschuivingen aan zijn randen, zoals Tibet vandaag doet.
\textbf{18.} $f = 600/3300 = 0.182$; $k = 2.303/11 = 0.209$ per miljoen jaar; $\tau = 1/(kf)
= \qty{26.3}{Myr}$.
\textbf{19.} $3.0\,e^{-30/26.3} = \qty{0.96}{km}$.
\textbf{20.} $26.3\ln(3.0/0.5) = \qty{47.1}{Myr}$ na \qty{410}{Ma}: op \qty{362.9}{Ma}.
\textbf{21.} De extra korst is $h/f$: $3.0/0.182 = \qty{16.5}{km}$ daalt tot $0.5/0.182 =
\qty{2.75}{km}$, dus wordt er \qty{13.75}{km} gesteente verwijderd.
\textbf{22.} $2700\times9.8\times13\,750 = \qty{3.6e8}{Pa}$, ongeveer \qty{0.36}{GPa}.
\textbf{23.} Ja: tegen het Laat-Devoon is de modelgordel gezakt tot ongeveer \qty{1}{km} gemiddelde
hoogte en daalt hij nog, met reliëf tussen zijn dalen en ruggen; rode grindlagen kunnen zich
lang vóór \qty{358.9}{Ma} in zijn laagten verzamelen, zoals bij Siccar Point, waar ze op een oppervlak
met reliëf rusten.
\textbf{24.} \textbf{\qty[text-series-to-math=true]{2880}{km}}, op \qty{500}{Ma}, toen
de subductie begon.
\end{solution}
